\chapter[Apresentação da Trajetória do Aluno]{APRESENTAÇÃO DA TRAJETÓRIA DO ALUNO}

Minha aproximação com tecnologia começou pela criação visual. Antes de pensar
em Engenharia de Software, explorava ferramentas como Photoshop, Premiere,
Blender e Animate para entender como imagens, vídeos e animações eram
construídos. Essa curiosidade despertou uma pergunta que se tornou decisiva na
minha trajetória: como funciona o software por trás dessas experiências? Durante
a quarentena, essa inquietação ganhou força. Tranquei o curso de Economia,
experimentei a criação de conteúdo e percebi que me motivava mais construir
soluções do que apenas comunicar ideias. Foi então que decidi iniciar minha
formação em Engenharia de Software.

Antes da AGES, o estágio de dois anos no Palácio da Justiça foi importante para
minha maturidade profissional. Embora a atuação estivesse baseada em ferramentas
\textit{no-code}, a experiência desenvolveu minha capacidade de organizar
entregas, compreender requisitos, comunicar-me com diferentes pessoas e atuar
em equipe. A participação em processos internos, inclusive em entrevistas de
candidatos para a área de TI, também ampliou minha percepção sobre cultura de
equipe e responsabilidade profissional. Posteriormente, a Dell IT Academy
aprofundou meu contato com arquitetura, qualidade e práticas modernas de
desenvolvimento, oferecendo uma base que se tornou muito relevante nas etapas
seguintes da graduação.

Minha primeira experiência na AGES I ocorreu no projeto Let Me Trial, uma
plataforma web voltada a apoiar médicos na identificação de pacientes elegíveis
para ensaios clínicos. Entrei no projeto sem conhecer a equipe e com muita
insegurança diante dos debates presenciais. Naquele período, questões pessoais
tornaram a participação nos encontros cada vez mais difícil. Ao perceber que
não conseguiria contribuir da forma que esperava, conversei com a liderança e
optei por trancar a disciplina. Foi uma decisão difícil, mas responsável: em
vez de permanecer no projeto sem condições de acompanhá-lo adequadamente,
escolhi me reorganizar para retornar em um momento melhor.

O retorno à AGES I, no projeto VIR, teve um significado diferente. Eu ainda não
dominava Git e não compreendia plenamente como um projeto se dividia entre
frontend, backend e banco de dados, mas já não tinha medo de fazer perguntas.
O tratamento da ansiedade, o conhecimento acumulado nesse intervalo e a
postura acolhedora dos AGES IV Luiz Felipe Porto e Arthur Maia contribuíram para
essa mudança. Eles criavam um ambiente leve, demonstravam disponibilidade para
explicar dúvidas desde as mais simples e valorizavam o bem-estar da equipe.
Nesse contexto, pude concentrar minha energia em aprender e desenvolver.

No VIR, comecei a reconhecer minha afinidade com o frontend de maneira mais
concreta. A implementação da \textit{footer}, com o roteamento de seus botões,
foi uma entrega marcante por representar uma funcionalidade completa sob minha
responsabilidade. Também busquei aprimorar a experiência visual para além do
requisito mínimo, criando, por exemplo, uma animação de rotação para o ícone de
configurações ao passar o cursor sobre ele. Pequenos detalhes como esse me
mostraram o quanto eu gostava de escrever código e ver a interface ganhar vida.
Mesmo que o projeto não tenha alcançado todas as entregas previstas, encerrei a
etapa satisfeito com minha contribuição individual e, principalmente, com a
confiança recuperada para aprender em público e colaborar com o time.

Na AGES II, no projeto Pró-Mata, entrei mais seguro para participar de
discussões, propor ideias e assumir decisões no frontend. A composição da equipe
concentrou muitos alunos da AGES II na frente de desenvolvimento, o que fez com
que minha atuação permanecesse voltada à construção do produto, sem uma
participação direta no banco de dados. Ainda assim, a experiência foi bastante
desafiadora: trabalhei com tecnologias que até então eram novas para mim, como
TanStack Router, TanStack Query, Zod, Zustand e Axios. A criação de fluxos no
Figma, a evolução do \textit{design system} e a implementação da versão
bilíngue da aplicação fortaleceram tanto meu domínio técnico quanto minha visão
de produto. Passei a entender que prototipar, discutir fluxos e antecipar
decisões visuais também são formas relevantes de contribuir para a engenharia
de software.

A AGES III, no FinUp, ampliou novamente essa perspectiva. Além de continuar
desenvolvendo funcionalidades para o aplicativo mobile, participei de decisões
sobre a organização dos repositórios, configuração de ambientes, revisões de
Pull Requests e integração entre frontend, backend, banco de dados, serviços em
nuvem e inteligência artificial. Ajudar a resolver problemas nos encontros
presenciais e nas chamadas de Discord, especialmente nas vésperas de
apresentações, foi um sinal concreto de que eu conseguia contribuir também fora
da minha zona de conforto. A etapa consolidou a percepção de que não preciso
saber tudo antes de participar de uma discussão técnica; preciso estar disposto
a investigar, pedir ajuda, testar hipóteses e validar soluções com o time.

Hoje, reconheço o frontend como a área em que tenho maior domínio e onde minha
inclinação criativa encontra espaço na construção de produtos. Ao mesmo tempo,
a trajetória na AGES tornou mais clara minha direção para me tornar um
profissional full stack. Quero aprofundar meus conhecimentos em banco de dados
e infraestrutura em nuvem, sem abandonar a preocupação com experiência de uso,
produto e qualidade visual. Levo das três etapas uma postura mais madura diante
de problemas, uma confiança construída pela prática e a convicção de que a
aprendizagem contínua e a colaboração são essenciais para construir soluções
reais.
